\lecture{7}{Thu 30 Oct 2025 14:01}{Gauge invariance}

So in QM a gauge transform is
\[
	\ket{\Psi }\to \ket{\widetilde{\Psi } }=\exp \left( \frac{iv}{\hbar} \Lambda (\mathbf{x} ,t) \right) \ket{\Psi }
.\]

So in classical E\&M, the Gauge potential $A$ is a vector quantity satisfying $\mathbf{B} =\nabla \times A$. I write non-bold $A$ because typing Abf inputs $\Ab f$. Also
\[
	\mathbf{E} =-\nabla \Phi - \frac{\partial A}{\partial t} 
.\]
The most fundamental law of E\&M is conservation of charge
\[
	\frac{\partial \rho }{\partial t} +\nabla \cdot \mathbf{J} =0
.\]
Gauge freedom says that physics does not change under transformations
\[
	A\to A+\nabla \Lambda 
\]
or
\[
	\Phi \to \Phi  -\frac{\partial \Lambda }{\partial t} 
\]
for some scalar $\Lambda $. Only potentials that can be connected by a scalar are valid. For some reason we write
\[
	\mathbf{E} =\frac{-\mu_0 c^2}{4\pi } \frac{e}{r^2} \hat{r} 
\]
instead of
\[
	-\frac{1}{4\pi \epsilon _0} \frac{e}{r^2} \hat{r} 
\]
for a point charge of charge $e$. If we do a very long discussion then we find the \emph{relativistic} energy for a charged particle in an EM field is given by swapping $\vec{p} $ for $\vec{p} -qA$.  Here $\vec{p} =-i\hbar \nabla $ and $A$ is a Gauge trasformation or something. SO this is the relativistic momentum operator or something.
